\section{Problem and Evaluation Object}
\label{sec:problem_v23}

An unobserved mode $h_t$ determines the transition kernel and vector reward
$\mathbf r_t\in\mathbb R^m$. Modes change at unknown times. The learner observes
states, actions and rewards, and receives a preference $w\in\Delta^{m-1}$.
Training uses discounted scalar utilities $w^\top\mathbf Q$ with $\gamma=0.99$.
The primary evaluation object is the final preference-conditioned policy,
which computes context from its own observed transition history. This policy
is history-dependent through its context even though its actor parameters are
fixed at evaluation. A retained set $A$ of such policies defines a separate
archive-coverage outcome.

For a locked mode $h$, let
$J_h(\pi,w)=\mathbb E[\sum_{t=0}^{L-1}\mathbf r_t]$, where an episode ends on
termination or at $L\leq1000$ steps. These reported returns are undiscounted.
For the fixed evaluation grid $W$, the empirical candidate set is
$F_h(A)=\{\widehat J_h(\pi,w):\pi\in A,w\in W\}$, using three episodes per
candidate. We compute $\mathrm{HV}_r(F_h(A))$ with a fixed reference $r$ and
summarize environments by $H^{-1}\sum_h\mathrm{HV}_r(F_h(A))$. Mode-conditioned
fronts are not pooled into a cross-mode Pareto front. In archive experiments,
the final-actor subset and full archive are evaluated from the same rollouts.

The primary outcome in the twenty-seed study is same-preference final-policy utility,
\begin{equation}
 U(\pi)=\frac{1}{H|W|}\sum_{h=1}^H\sum_{w\in W}
              w^\top\widehat J_h(\pi,w).
 \label{eq:direct_utility}
\end{equation}
Here the sum runs over all modes and grid weights; there is no maximization
over other preferences or archived actors. A policy-set HV gain need not
increase this deployment-oriented outcome.

\paragraph{Why utility and coverage differ.}
Consider one mode and two requested weights $e_1=(1,0)$ and $e_2=(0,1)$.
A policy returning $(2,1)$ at $e_1$ and $(1,2)$ at $e_2$ has direct utility
$2$. A second policy that swaps these two responses has direct utility $1$.
Their return sets, and hence HV at any common reference, are identical.
This illustrative example separates preference-conditioned execution from
set coverage without assuming either metric determines the other.

\section{BAMOR Update Portfolio}
\label{sec:method_v23}

BAMOR denotes the frozen v25 implementation throughout the current comparisons.
At each interaction it constructs context from past transitions and conditions
one shared actor and twin vector critics on that context and a preference.
At each block boundary it chooses a fixed update configuration from a portfolio.
The configuration governs both collected preferences and replay updates for
the block; it does not instantiate a separate actor. Training rewards update
the selector, while test returns and mode identities are reserved for
evaluation. Algorithm~\ref{alg:bamor} gives the complete sequence.

\subsection{Context and replay}

One actor $\pi_\theta(a\mid s,w,z)$ and twin vector critics
$\mathbf Q_{\phi_i}(s,a,w,z)$ share a history-derived context $z$: reward
exponential averages, a scalar coordinate-response statistic and an
eight-dimensional learned transition embedding. The encoder reconstructs
scaled rewards and state displacements with temporal contrastive regularization.
Replay stores the collected context; preference relabeling changes only its
weight coordinates, not its stored history features. Mode labels are not
inputs. Appendix~\ref{app:implementation_details} specifies the coordinate
feature, encoder loss and exponential averages.

\subsection{Update configurations}
\label{sec:operator_portfolio}

Each arm $o$ specifies a preference sampler, relabel rate and copy count $C_o$,
critic candidate count $K_o^Q$, actor candidate count $K_o^\pi$, vector-loss
coefficient $\lambda_o$, and behavior-cloning coefficient $b_o$. There are seven
coverage arms, seven recovery arms and four extreme-preference arms, comprising
14 distinct sampler/update settings. Identical settings in different arm slots
retain separate scores and allocation counts. Relabeling
reuses vector-reward transitions under new preference inputs. All arms update
the same networks; the family names describe their intended roles. The full
configuration and warmup sequence are in Appendix~\ref{app:v23_operators}.

\paragraph{Critic update.}
For each replay transition, let $u$ denote its update preference after relabeling.
Set candidate preference $v_1=u$ and sample the remaining preferences using the
arm's replay sampler. Draw $a'_j\sim\pi_\theta(\cdot\mid s',v_j,z')$ and define
\begin{align}
 \mathbf q'_j&=\min_{i=1,2}\mathbf Q_{\phi_i^-}(s',a'_j,v_j,z'), \nonumber\\
 j^*&=\arg\max_j\{u^\top\mathbf q'_j-\alpha\log\pi_\theta(a'_j\mid s',v_j,z')\},\nonumber\\
 \mathbf y_v&=\mathbf r+\gamma(1-d_{\mathrm{term}})
       [\mathbf q'_{j^*}-\alpha\log\pi_\theta(a'_{j^*}\mid s',v_{j^*},z')\mathbf 1],\nonumber\\
 y_s&=u^\top\mathbf y_v.
 \label{eq:v25_targets}
\end{align}
Here the minimum is componentwise, targets are detached, $\alpha=0.2$ is fixed,
and only true termination sets $d_{\mathrm{term}}=1$; time limits retain the
bootstrap. The critic loss is
\begin{equation}
 \mathcal L_Q=\mathbb E\sum_{i=1}^2\left[
 (u^\top\mathbf Q_{\phi_i}(s,a,u,z)-y_s)^2+
 \frac{\lambda_o}{m}\|\mathbf Q_{\phi_i}(s,a,u,z)-\mathbf y_v\|_2^2\right].
 \label{eq:v23_critic_loss}
\end{equation}
Both target terms include the same selected-candidate entropy bonus.

\paragraph{Actor update.}
Generate $K_o^\pi$ actions $a_j\sim\pi_\theta(\cdot\mid s,v_j,z)$ with $v_1=u$.
Unlike the critic target above, rank these actions with critics conditioned on
the update preference $u$:
\begin{align}
 \mathbf q_j&=\min_{i=1,2}\mathbf Q_{\phi_i}(s,a_j,u,z),
 &j^\dagger&=\arg\max_j u^\top\mathbf q_j,\\
 \mathcal L_\pi&=\mathbb E\left[
 \alpha\log\pi_\theta(a_{j^\dagger}\mid s,v_{j^\dagger},z)
 -u^\top\mathbf q_{j^\dagger}
 +\frac{b_o}{d_a}\|a_{j^\dagger}-a\|_2^2\right].
 \label{eq:v23_actor_loss}
\end{align}
The discrete ranking is detached; gradients flow through the selected sampled
action and its generating policy log probability. $d_a$ is action dimension and
$a$ in the last term is the replay action. Thus multi-candidate updates are
specified training heuristics, not an exact optimality operator.

\subsection{Progress-driven block allocation}
\label{sec:method_v25}

Progress is measured in reward coordinates without multiplying by each
transition's collection preference. The first 5,000 transitions initialize
causal reward moments. Before each subsequent 1,500-transition block $k$, freeze
the per-objective standard deviations $s_{k,j}$ from past rewards. For the
first- and second-half means, the selected arm receives
\begin{equation}
 P_k=\frac{1}{m}\sum_{j=1}^m\tanh\left(
 \frac{\overline r^{(2)}_{k,j}-\overline r^{(1)}_{k,j}}{s_{k,j}}\right),
 \qquad q_o\leftarrow0.87q_o+0.13P_k,
 \label{eq:v25_progress}
\end{equation}
The arm stays fixed within a block, while interaction preferences change at
each transition. A completed block updates its score after the last reward
and before that transition's gradient; incomplete blocks remain unscored.
After ten prescribed warmup blocks, the selector first corrects family
allocation deficits against fractions $(0.36,0.22,0.18)$. Otherwise it samples
a family using the mean of its two highest arm scores, then an arm using its
own score, with softmax and uniform exploration at both levels. The score-blind
control substitutes zero routing scores while computing the same observations.
Exact temperatures, scales and scheduling parameters are in
Appendix~\ref{app:implementation_details}.

\subsection{Implementation budget and optional archive}

The main protocol uses 75,000 interactions and evaluates only the final actor
with its final encoder, without intermediate evaluation or actor selection.
Training context trackers persist across resets; evaluation initializes them
afresh. The separate 50k protocol retains 19 actors at steps
$5{,}000,7{,}500,\ldots,50{,}000$, all evaluated with the final encoder.
Archive evaluation does not feed the selector. Per-run network, optimization
and candidate resources are in Table~\ref{tab:v23_budget}.

\begin{algorithm}[t]
\caption{BAMOR-v25 training with optional actor retention}
\label{alg:bamor}
\label{sec:algorithm}
\begin{algorithmic}
\STATE Initialize networks, replay, context trackers and arm scores $q_o=0$; set $L=5{,}000$
\FOR{$t=1,\ldots,T$}
 \IF{$t\leq L$}
  \STATE Use the first warmup arm
 \ELSIF{a new scored block starts}
  \STATE Select an arm by Section~\ref{sec:method_v25}; freeze scales from past rewards
 \ENDIF
 \STATE Sample the arm's interaction preference; construct context from past transitions
 \STATE Act randomly if $t<L$, otherwise use the actor; observe the transition
 \STATE Update reward moments; for $t>L$, accumulate half-block rewards and allocation counts
 \IF{the 1,500-transition scored block is complete}
  \STATE Update $q_o$ by Eq.~\eqref{eq:v25_progress}, increment its block count and close the block
 \ENDIF
 \STATE Update context trackers; store the transition, both contexts and true-termination mask
 \IF{$t\geq L$}
  \STATE Relabel and expand a replay minibatch according to the selected arm
  \STATE Update critics using Eq.~\eqref{eq:v25_targets}, actor and target critics; update encoder
 \ENDIF
 \STATE Retain actor parameters at the prescribed steps, if an archive is requested
\ENDFOR
\STATE Leave any incomplete block unscored; evaluate the final actor in each locked mode
\STATE For archive experiments, also evaluate saved actors with the final encoder
\end{algorithmic}
\end{algorithm}

\FloatBarrier
\section{Supporting Mathematical Properties}

\paragraph{Target consistency and feedback.}
Simplex weights project the soft vector target exactly onto its scalar target
in Eq.~\eqref{eq:v25_targets} (Proposition~\ref{prop:soft_projection}). The
progress statistic and its score EMA stay in $[-1,1]$, and constant reward
offsets cancel at fixed scale (Proposition~\ref{prop:progress}). These are
properties of the implemented quantities.

\paragraph{Candidate selection and local stability.}
At equal weights, perturbing $(1,0),(0,0.9)$ to $(1,0),(0,1.1)$ changes the
unique selected vector's first coordinate by $1$ despite an input change of
only $0.2$ (Proposition~\ref{prop:reselection}). Scalar-maximization stability
therefore does not imply selected-vector stability. For matched actions and
contexts, let $S_j=u^\top q_j+e_j$, with critic-coordinate perturbations at
most $\epsilon\geq0$ and entropy-bonus perturbations at most $\rho\geq0$. If the winner
$j^*$ satisfies
\begin{equation}
 S_{j^*}-S_j>2\delta\quad(j\neq j^*),\qquad\delta=\epsilon+\rho,
 \label{eq:main_candidate_margin}
\end{equation}
the winner is unchanged. With fixed reward and bootstrap multiplier $b\geq0$,
the soft targets obey $\|Y-\widetilde Y\|_\infty\leq b\delta$
(Proposition~\ref{prop:candidate_margin}).

Appendix~\ref{sec:theory} gives the Lean proofs and conditional tracking and
archive results. Tracking requires contractive selected maps, bounded update
error and bounded selected-map residual. HV monotonicity fixes returns,
reference and encoder; archive selection additionally requires uniform return
accuracy. These assumptions delimit the supporting analysis.
